\noindent {\fontsize{13 pt}{12 pt}\selectfont{\textbf{Background and Research Interests -}}}

\vspace{0.10 cm}

My Masters in Physics at \textbf{National Institute of Technology Jamshedpur (NIT)} sparked my interest in Astrophysics and Computation, so for my Master’s thesis I joined \textbf{Indian Institute of Astrophysics (IIA)} and worked on ‘Optimisation of feature tracking code to analyse small-scale magnetic features on the solar photosphere’. I redeveloped this code enhancing its runtime from 7 hrs to 2.5 hrs in Python and from 1 hr to 25 min} in Julia outperforming the IDL code with runtime of 40 min written by my supervisor Dr. Anusha \cite{2017A&A...598A..47A}. This code is now being used to analyse SUNRISE-III data and will soon become open-source, resulting in publication. I then joined \textbf{Inter-University Center for Astronomy and Astrophysics (IUCAA)} and worked on the analytical and numerical modelling of compact stars, using two-fluid model and fermion-boson star \cite{Diedrichs:2023trk} to study the fraction of dark matter inside their core and it's effect on physical parameters (M-R, tidal deformability, compactness etc.) 

\vspace{0.10 cm}

Currently, I am working with \textbf{Dr. Jagdish C. Joshi} at \textbf{Aryabhatta Research Institute of Observational Sciences (ARIES), India} on Synchrotron-SSC model of Gamma Ray Burst (GRB). We studied minimum variability time scales (MVT), parameter degeneracy in GRB models using Markov Chain Monte Carlo algorithms (Bilby, emcee, PyMultinest) and non-thermal very high energy (VHE) emissions through radiative processes like Synchrotron, Inverse Compton, Synchrotron Self Compton (SSC) during prompt and afterglow phase of GRBs. We worked out analytical calculations for Synchrotron-SSC model \cite{Beniamini_2013}. Supervision of Dr. Jagdish has been very effective, he introduced me to GRBs, Pulsar Wind Nebulae (PWNe), Supernova Remnants (SNRs) and astronomy using Cosmic Rays (CRs), $\gamma$-rays and neutrinos. In this project I learned analytical modeling of non-thermal emission mechanisms and high-energy (HE) astrophysical data analysis, which aligns with my ambition to study origin of CRs and particle acceleration in SNRs, PWNe, etc. During this period, I was introduced to research of \textbf{Dr. Martina Cardillo} and got an opportunity to interact with her which sparked my interest in \textbf{ASTRI Mini-Array}.

\vspace{0.3cm}
\noindent {\fontsize{13 pt}{12 pt}\selectfont{\textbf{Why PhD in AASS?}}}

\vspace{0.10 cm}

The PhD at the \textbf{Università degli Studi di Roma Tor Vergata} with \textbf{Istituto di Astrofisica e Planetologia Spaziali (INAF-IAPS)} represents a valuable opportunity to advance my research in HE astrophysics using CR and neutrinos. Dr. Martina’s expertise lies in $\gamma$-ray and CRs, who is also a key figure in the \textbf{ASTRI} project, and her leadership in this field perfectly align with my research goals. The active collaboration between the \textbf{ARIES}, \textbf{ASTRI} and \textbf{IceCube} create an ideal environment for me to pursue high-impact research and continue my passion for HE astrophysics.


\vspace{0.3cm}
\noindent {\fontsize{13 pt}{12 pt}\selectfont{\textbf{Research Goals -}}}

\vspace{0.10 cm}
\textbf{Dr. Martina}'s project \textbf{Gamma-ray Astronomy and the Search for Lost Pevatrons (with a Little Help from Neutrinos)} provides a robust framework for investigating galactic CR for which I aim to follow strict project timelines over the next three years, leveraging my computational and analytical expertise to drive meaningful progress.

\vspace{0.10 cm}
Let me briefly summarise my research objectives. I aim to start with a literature study focusing on CR origin and existing $\gamma$-ray emission models for SNRs with Dr. Martina's team. I plan to leverage my experience with the $\gamma \Pi$ Python package to analyze VHE observations from instruments such as \textbf{ASTRI, LST-1} and \textbf{CTAO}. Additionally, with my expertise developed at IIA, I can contribute to Dr. Martina's team to convert and optimize their existing FORTRAN routines into efficient Python code. My programming skills will help me to adapt and contribute to ongoing efforts of her team very quickly. After this, I will work on refining and advancing the SNR acceleration model, incorporating neutrino emission collaborating with my current guide \textbf{Dr. Jagdish} at \textbf{ARIES} with his expertise in astrophysical neutrinos to differentiate between leptonic and hadronic processes. This will allow us to compare models across source types and apply our theoretical and computational frameworks to simulated/real existing data from current Cherenkov Telescopes to test the model, helping us in the multiwavelength and multimessenger approach. After testing these models, we will eventually apply them to study VHE and neutrino data from \textbf{ASTRI} and \textbf{IceCube}, aiming to identify PeVatrons and plan our future observational strategies.

\vspace{0.10 cm}

My background in modeling GRBs and parameter estimation equip me to excel in this project. My skills in code optimisation, numerical and statistical frameworks, efficient pipeline developments, data analysis, collaborative teamwork, scientific writing, compatibility with Linux Machines (servers, HPC (High Performance Computing) and Clusters), handling large datasets, and presentation and seminars will help me to contribute to Dr. Martina's group effectively.

\vspace{0.3cm}
\noindent {\fontsize{13 pt}{12 pt}\selectfont{\textbf{Diversity and Outreach -}}}

\vspace{0.10 cm}
I am also committed to enriching the scientific community through diversity and outreach activities. In collaborative and international settings, I value diverse perspectives from both diverse research and cultural backgrounds. Scientific communication and criticism further enhance our studies in the right direction. I am passionate about scientific communication through conferences, seminars, workshops and public talks which will be an additional asset for educational and outreach programs at the university and INAF-IAPS, sharing the exciting research in HE astrophysics with broader audiences. In particular, Dr. Martina’s deep engagement in outreach activities and her proactive leadership in diversity initiatives create an ideal platform to which I can contribute effectively.

\vspace{0.3cm}

This PhD opportunity with \textbf{Dr. Martina} and collaboration with \textbf{Dr. Jagdish} will allow me to grow as an independent researcher and excel in this field, and I am committed to delivering my best.



